\chapter{重要度依据总表（文件证据）}
\label{ap:evidence}

\begin{tishi}
本附录是全书\textbf{星级判定的原始依据}。所有星级都不是主观印象，而是从下列 6 份文件里读出来的：
每一条都注明\textbf{证据编号}与\textbf{原文摘录}，可以逐条回查。
\end{tishi}

\section{证据文件清单 E1--E6}

\subsection*{E1\quad《8-1 考情分析》PPT（能动考研笑开研·小刘学长，西华大学 24 级能动专硕）}

\noindent\textbf{第 15 页——题型与分值（原文照录）：}

\begin{quote}\small
\textbf{2018 年以前：}\\
选择题 30 分（10 道题，每题 3 分）\\
填空题 30 分（10 道题，每题 3 分）\\
计算题 90 分（6--7 道题，每题 10--15 分）\\
作图题（只考过一次，7 分）\\
证明题（变相的计算题）12--15 分\\
判断题 20 分（10 道题，每题 2 分）\\[4pt]
\textbf{2019 年以后开始：}\\
选择题 30 分（10 道题，每题 3 分）\\
填空题 30 分（10 道题，每题 3 分）\\
判断题 20 分（10 道题，每题 2 分）\\
简答题 12 分（3 道题，每题 4 分）\\
计算题 58 分（3 道题，每题 15--20 分）
\end{quote}

\noindent\textbf{第 16 页——各题型考法（原文照录）：}

\begin{quote}\small
选择判断：分析概念、混淆概念的方法等\\
填空题：名称解释的定义、小型计算题等\\
简答题：水头损失的分类及定义、均匀流的特点等\\
计算题：伯努利方程、动量方程等
\end{quote}

\noindent\textbf{第 25 页——出题风格（原文照录）：}

\begin{quote}\small
① 重视基础，新题型很少\\
② \textbf{真题重复率很高}\\
③ 题量适中（考试一个半小时基本就能学完）\\
④ 计算量适中\\
⑤ 围绕参考书（大纲）进行复习即可\\
⑥ 总体难度中等（把握细节可冲高分）
\end{quote}

\noindent\textbf{第 33 页——逐章重要度（原文照录，含原文笔误）：}

\begin{quote}\small
第一章：绪论\qquad\qquad\qquad\qquad\ 重点掌握基础知识概念\\
第二章：流体静力学\qquad\qquad\qquad 计算题基础入门，重点掌握方法\\
第三章：流体运动学\qquad\qquad\qquad 公式较多，重点记忆\\
第四章：流体动力学基础\qquad\qquad \textbf{考试计算题重点}\\
第五章：管口、孔口、管嘴的水力计算\quad\ \textbf{考试计算题重点}\\
第六章：相似理论与量纲分析\qquad\quad 量纲之间的转化、相似理论\\
第七章：理想流体动力学\qquad\qquad\quad 概念为主\\
第八章：黏性流体动力学基础\qquad\quad 会涉及选择题和填空题考点
\end{quote}

\begin{yicuo}
第 33 页把第五章写成"管口、孔口、管嘴的水力计算"，教材原目录是\textbf{"管路、孔口、管嘴的水力计算"}
（"管口"应为"管路"），属 PPT 录入笔误，此处照录以存真。
\end{yicuo}

\subsection*{E2\quad《9-1 考点分析》PPT（同一作者）}

\noindent\textbf{第 4 页——分值分布与章节分层（原文照录）：}

\begin{quote}\small
西华大学的专业课题型主要包括选择题、填空题、判断题、简答题、计算题。
\textbf{计算题占考试主要分数（50\%）}，其次是选择题和填空题。
\textbf{第一章到第三章}考点涉及概念性的知识点，掌握基本计算公式和知识解释即可。
\textbf{第四章}为考试主要考点，需掌握题型的计算套路。
\end{quote}

\noindent\textbf{第 5 页（原文照录）：}

\begin{quote}\small
\textbf{第五章}为考试的主要计算考点，同第四章一样需掌握计算套路。
\textbf{第六章到第八章}，少许的填空、判断、选择题会涉及；
主要掌握学习基本知识点概念、计算公式、学会套用计算公式即可。
\end{quote}

\subsection*{E3\quad《西华大学 818 工程流体力学考研分析》Word 文档}

\noindent\textbf{重点段落（原文照录）：}

\begin{quote}\small
针对初试参考书本，市面上有很多关于工程流体力学的版本，在这里个人推荐官网书籍
《流体力学与流体机械》--赵琴，初试的真题也是根据这本书来进行。
参考书共有 \textbf{14 个章节，初试考察的知识点只包含前 8 章}，总的来说把前 8 章学扎实高分不是问题；
也没必要再增加一些不重要知识点来学习，反而增加多余的负担。\\
参考书 \textbf{9--14 章的内容应对复试}，结合复试参考书《流体机械原理》进行学习……\\[4pt]
\textbf{每年的真题大部分都源于课后习题的改编，课后习题一定要特别认真对待。}
增加历年常考选择题型重点带过。\\[4pt]
学习资料：1）工程流体力学知识点速过\quad 2）考点精讲及复习思路视频\quad 3）笔记汇总\quad
4）流体力学-教案 PPT-赵琴\quad 5）\textbf{历年真题（16--26）、课后习题、题库}
\end{quote}

\subsection*{E4\quad 历年真题（2014--2022 年原卷 + 2026 年回忆版）}

\noindent 每道真题都能定位到\textbf{年份 + 题型 + 题号}。真题是星级判定中最硬的一条证据：
某个考点在真题里\textbf{出现的次数越多、位置越靠后（计算题）}，星级越高。
本附录第 \ref{sec:apA-tab} 节给出逐章对应表。

\subsection*{E5\quad 题库（《流体力学与流体机械》试题库（四）、西华选择题题库精简版）}

\noindent 题库题目与真题高度重合（E1 第 25 页"真题重复率很高"、2026 回忆版"和往年真题还有题库类似"），
是判断一个概念"会不会考、以什么形式考"的第二类证据。

\subsection*{E6\quad 赵琴《流体力学与流体机械》教材（第 1--8 章）}

\noindent 教材决定\textbf{考纲范围}：第 1--8 章内的知识点才可能是考点，第 9--14 章（气体动力学、流体机械）
本次初试不考。教材还决定\textbf{公式形式与符号}（如 $h_w$、$\alpha$、$\beta$ 的取值），
答题时按教材符号书写最稳妥。

\section{逐章考点星级与依据}
\label{sec:apA-tab}

{\small
\begin{longtable}{@{}p{0.13\textwidth}p{0.30\textwidth}cp{0.44\textwidth}@{}}
\toprule
章 & 考点 & 星级 & 主要依据 \\
\midrule
\endfirsthead
\toprule
章 & 考点 & 星级 & 主要依据 \\
\midrule
\endhead
\bottomrule
\endlastfoot

\multicolumn{4}{@{}l}{\textbf{第 1 章\quad 绪论}}\\
1 & 连续介质假设 & ★★★★ & E4 2022 单选 1；E5 题库选择 \\
1 & 黏性与牛顿内摩擦定律 & ★★★★★ & E4 2026 计算 1（牛顿流体斜板）；E5 题库计算 \\
1 & 牛顿流体与非牛顿流体 & ★★★★ & E5 题库选择 2（汽油＝牛顿流体）\\
1 & 密度、比容、重度 & ★★★ & E4 2015 填空 1（851\,kg/m$^3$ 求重度）\\
1 & 压缩性与膨胀性 & ★★★ & E5 题库选择；E6 教材 1.3.2 \\
1 & 质量力与表面力 & ★★★ & E5 题库选择 3（体积力有势）\\
1 & 表面张力 & ★★ & 教材小节，真题未出现 \\
1 & 流体力学研究任务与方法 & ★★ & 概念题素材 \\
\midrule
\multicolumn{4}{@{}l}{\textbf{第 2 章\quad 流体静力学}}\\
2 & 静力学基本方程与测压管水头 & ★★★★★ & E4 2022 单选 3、单选 7；E1 slide33"计算题基础入门"\\
2 & 平面上的总压力 & ★★★★★ & E4 2015 计算（圆柱闸门）；E1 slide33 \\
2 & 曲面上的总压力 & ★★★★★ & E5 题库计算 1（弧形闸门 $F_x,F_z$ 与力矩）\\
2 & 压力体 & ★★★★★ & 同上；E4 各年作图/计算 \\
2 & 浮力与阿基米德原理 & ★★★★★ & E1 slide33；E4 历年计算 \\
2 & 欧拉平衡微分方程 & ★★★★ & E4 2022 单选 5（$dp$ 形式）；E5 题库选择 3（有势）\\
2 & 等压面 & ★★★★ & E5 题库；E6 教材 2.2 \\
2 & 压强表示法（绝对/相对/真空） & ★★★★ & E4 2022 单选 7（虹吸管压强）；E5 题库选择 7\\
2 & 静压强测量（测压计） & ★★★★ & E4 2015 填空（U 形水银压差计）\\
2 & 流体静压强两大特性 & ★★★★ & E2、E5 概念题高频\\
2 & 压强分布图 & ★★★★ & E4 早期作图题（7 分，仅一次）\\
2 & 等加速直线运动相对平衡 & ★★★ & E6 教材 2.4；E5 题库\\
2 & 等角速度旋转相对平衡 & ★★★ & E6 教材 2.4；E5 题库\\
2 & 惯性矩 $J_C$ 与形心 & ★★★ & 平面总压力计算配套 \\
\midrule
\multicolumn{4}{@{}l}{\textbf{第 3 章\quad 流体运动学}}\\
3 & 质点导数与加速度 & ★★★★★ & E4 各年选择/填空；E1 slide33"公式较多，重点记忆"\\
3 & 流线与迹线 & ★★★★★ & E4 2015 计算 1（流线方程、流动方向）\\
3 & 流线方程 & ★★★★★ & 同上；E5 题库 \\
3 & 连续性方程 & ★★★★★ & E4 2022 单选 4（方管进压缩机）；E5 题库 \\
3 & 过流断面、流量、平均流速 & ★★★★ & E4 2022 单选 4；E5 题库\\
3 & 系统与控制体 & ★★★★ & 伯努利/动量方程推导基础 \\
3 & 旋转角速度与有旋判据 & ★★★★ & E5 题库选择 13（涡量为零）\\
3 & 拉格朗日法与欧拉法 & ★★★ & E1 slide33 \\
3 & 定常/非定常、均匀/非均匀 & ★★★ & E1 slide16 简答（均匀流特点）\\
3 & 一元、二元、三元流动 & ★★★ & 概念题 \\
3 & 脉线、流管、元流 & ★★★ & 概念题 \\
3 & 湿周与水力半径 & ★★★ & 明渠与管路计算配套 \\
3 & 流体微团运动分解 & ★★★ & E6 教材 3.4 \\
3 & 线变形率与角变形率 & ★★★ & E6 教材 3.4 \\
\midrule
\multicolumn{4}{@{}l}{\textbf{第 4 章\quad 流体动力学基本方程（E1/E2 点名"考试计算题重点"）}}\\
4 & 实际流体总流的伯努利方程 & ★★★★★ & E1 slide16"计算题：伯努利方程"；E1/E2 slide33/5"第四章＝考试计算题重点"\\
4 & 动量方程及其分量式 & ★★★★★ & E4 2015 计算（水平弯管求 $F_x,F_y$）；E4 2026 计算 3（教材例题 4.12）；E1 slide16\\
4 & 弯管、射流、喷嘴类综合题 & ★★★★★ & 同上 \\
4 & 典型应用（文丘里、皮托、虹吸、孔口管嘴） & ★★★★★ & E4 2022 单选 6；E4 2026 计算 2（文丘里流量证明）\\
4 & 带泵的伯努利方程与泵功率 & ★★★★ & E5 题库计算 4（水泵工况点与节流调节）\\
4 & 理想流体运动微分方程（欧拉方程） & ★★★ & E6 教材 4.1 \\
4 & 兰姆-葛罗米柯方程 & ★★★ & E6 教材 4.1 \\
4 & 伯努利积分（理想流体沿流线） & ★★★ & E6 教材 4.2 \\
4 & 动量矩方程 & ★★ & E6 教材 4.4 \\
\midrule
\multicolumn{4}{@{}l}{\textbf{第 5 章\quad 管路、孔口、管嘴的水力计算（E1/E2 点名"主要计算考点"）}}\\
5 & 沿程水头损失与达西公式 $\lambda$ & ★★★★★ & E4 2022 单选 3、单选 10；E1 slide16 简答"水头损失分类及定义"\\
5 & 孔口出流 & ★★★★★ & E1 slide33；E2 slide5"第五章为考试的主要计算考点"\\
5 & 管嘴出流与圆柱形外管嘴 & ★★★★★ & E4 2022 单选 8（$H_0\le 9$ m、$l=(3\sim4)d$）；E4 2026 计算 5（教材管嘴例题）\\
5 & 简单长管水力计算 & ★★★★★ & E1 slide33；E5 题库 \\
5 & 复杂管路 & ★★★★ & E1 slide33；E6 教材 5.7 \\
5 & 短管、虹吸管、吸水管 & ★★★★ & E4 2022 单选 7（虹吸管）；E5 题库 \\
5 & 圆管层流（速度分布、流量、切应力） & ★★★★ & E4 2015 填空（平均速度＝0.5 倍最大速度）；E4 2015 选择（切应力分布）\\
5 & 雷诺实验与流态判别 & ★★★★★ & E4 2022 单选 9（两圆管流态）；E4 2026 简答 3\\
5 & 局部水头损失 & ★★★★ & E4 2015 计算（突然缩小）；E1 slide16\\
5 & 流动阻力与水头损失分类 & ★★★★★ & E1 slide16 简答"水头损失的分类及定义"；E1 slide33；E2\\
5 & 湍流与沿程阻力系数 & ★★★ & E6 教材 5.4 \\
5 & 水击（水锤） & ★★★★ & E4 2026 简答 1（水击现象）；E6 教材 5.8\\
\midrule
\multicolumn{4}{@{}l}{\textbf{第 6 章\quad 相似理论与量纲分析（E4 九年内 13 处出题，比 E2 的估计更重）}}\\
6 & 力学相似的三个层次（几何/运动/动力，$\lambda_l,\lambda_v,\lambda_t,\lambda_F$） & ★★★ & E4 2020 判断 8（动力相似是流动相似的主导因素）；E1 slide33\\
6 & 相似三定理（正定理、逆定理、$\pi$ 定理） & ★★★ & E1 slide33"第六章 量纲转化、相似理论"；E2"掌握基本概念即可"\\
6 & 三大相似准则 $\Rey$、$\Eu$、$\Fr$ & ★★★★ & E4 2014 选择 10（有压管流动力相似$\to$雷诺准则）；E4 2016 填空 4、E4 2022 单选 9（管道阀门实验）；E4 2019 判断 2（$\Fr$ 物理意义）；E5 题库选择 12\\
6 & 其他准则 $\mathrm{Ma}$、$\mathrm{St}$、$\mathrm{We}$、$\mathrm{Ca}$ & ★★★ & E4 2020 判断 9（可压缩流要求柯西数相等）；E1 slide16\\
6 & 定性/非定性准则与相似准则选取原则 & ★★★ & E4 2021 简答 3（有压管流与明渠分别选哪个准则）\\
6 & 量纲与单位、导出量纲速查 & ★★★★ & E4 2014 选择 9（切应力 $\tau$ 的量纲）；E4 2020 判断 10；E5 题库选择 11（运动黏滞系数 $\mathrm{L^2T^{-1}}$）\\
6 & 量纲和谐原理（齐次性） & ★★★★ & E1 slide33（量纲转化）；E2（公式套用）\\
6 & 瑞利法（含 $n\le4\sim5$ 限制） & ★★★★ & E4 2019 判断 8（瑞利法待求量纲指数不超过 4 个）；E1 slide33\\
6 & $\pi$ 定理／布金汉定理（六步套路） & ★★★★ & E4 2016 填空 3、E4 2018 选择 5、E4 2019 选择 5；E4 2020 选择 7（水泵功率 $N$）；E1 slide33\\
6 & 典型例题一：圆管沿程压强损失 $\Delta p$ & ★★★★ & E2（公式套用）；E3（真题多源于课后习题改编）\\
6 & 典型例题二：绕流阻力 $F$ 与孔口流量 & ★★★★ & E1 slide16（填空考小计算）；E2（公式套用）\\
6 & 泵与风机的相似换算（$Q/H/P$ 与 $n,D$） & ★★ & E5 题库选择 15（风机并联风量）；E4 2020 选择 7；E1/E2 未列为初试重点\\
6 & 比转速 $n_s$ 与标准状态换算 & ★★ & E5 题库问答 2（泵与风机标准条件与相似换算）\\
\midrule
\multicolumn{4}{@{}l}{\textbf{第 7 章\quad 理想流体动力学（E1 slide33"概念为主"）}}\\
7 & 势函数与流函数的存在条件、相互关系 & ★★★★★ & E4 2015 计算（流函数求速度与流线）；E4 2026 计算 4（证明势函数存在性）\\
7 & 基本平面势流（均匀流、源汇、偶极子、环流） & ★★★★ & E4 2015 计算 8（均匀流＋点源叠加求驻点）\\
7 & 势流叠加与驻点、圆柱绕流 & ★★★★ & 同上；E2 slide5 \\
7 & 旋涡运动与速度环量 & ★★★ & E1 slide33"概念为主" \\
7 & 旋涡基本定理（斯托克斯、汤姆孙、亥姆霍兹） & ★★★ & 概念题 \\
7 & 毕奥-萨伐尔定理与诱导速度 & ★★ & 概念题 \\
7 & 空间势流 & ★ & 扩展内容 \\
\midrule
\multicolumn{4}{@{}l}{\textbf{第 8 章\quad 黏性流体动力学基础（E1 slide33"会涉及选择题和填空题考点"）}}\\
8 & 边界层概念（厚度、两种流态、压强分布） & ★★★ & E1 slide33；E2 slide5 \\
8 & 边界层方程、量级比较与普朗特方程 & ★★★ & 同上 \\
8 & 边界层分离与形状阻力 & ★★★ & 同上 \\
8 & 平板边界层计算（层流/湍流/混合） & ★★ & E4 2015 计算 9（平板摩擦阻力比 $F_1/F_2$）\\
8 & 纳维-斯托克斯方程 & ★★★ & E1 slide33 \\
8 & N-S 方程的几种解析解 & ★★ & 扩展内容 \\
8 & 雷诺时均方程与湍流模型 & ★ & 扩展内容 \\
\end{longtable}
}

\section{真题年份--章节分布（据 2014--2022 原卷与 2026 回忆版整理）}

\begin{center}
\begin{tabularx}{\textwidth}{@{}lXXXXXX@{}}
\toprule
年份 & 选择 & 填空 & 判断 & 简答 & 计算 & 备注 \\
\midrule
2014 & 10 & 10 & --- & --- & 11 题选做 & 含明渠、流体机械 \\
2015 & 15 & 10 & --- & --- & 10 题选做 & 含明渠、流体机械 \\
2016 & 10 & 10 & --- & --- & --- & \\
2017 & 10 & 10 & --- & --- & --- & \\
2018 & 10 & 10 & --- & --- & --- & \\
2019 & 10 & 10 & 10 & 3 & 3 & 题型改革年 \\
2020 & 10 & 10 & 10 & 3 & 3 & \\
2021 & 10 & 10 & 10 & 3 & 3 & \\
2022 & 10 & 10 & 10 & 3 & 3 & \\
2026 & 10 & 10 & 10 & 3 & 5 & 回忆版 \\
\bottomrule
\end{tabularx}
\end{center}

\noindent\textbf{读法：}2019 年是分水岭——之前是"选择 30 ＋ 填空 30 ＋ 计算 90"，
之后固定为"选择 30 ＋ 填空 30 ＋ 判断 20 ＋ 简答 12 ＋ 计算 58"。
因此\textbf{2019 年之后的试卷结构才是你的应试模板}；做 2018 年以前的卷子时，
计算题数量多、分值高，但考点分布仍然有效。
